\begin{studiumworked}{M27. Ascenso vertical elemental y duración · Intermedio}
\textbf{Enunciado.} Un modelo vertical inicia con masa $427\,\mathrm{kg}$, consume 2 kg/s durante $55.4\,\mathrm{s}$ y tiene $c=2450$ m/s. Considere gravedad constante $g=9.80665$ m/s$^2$, sin resistencia. Halle empuje, masa final, incremento ideal, contribución gravitatoria e incremento de velocidad. Compruebe empuje-peso inicial.

\textbf{Datos.} Masa inicial: $427\,\mathrm{kg}$; Consumo: $2\,\mathrm{kg/s}$; Duración: $55.4\,\mathrm{s}$; Velocidad efectiva: $2450\,\mathrm{m/s}$; Gravedad del modelo: $9.80665\,\mathrm{m/s^2}$.

\textbf{Incógnitas.} Incremento real del modelo y viabilidad inicial.

\textbf{Hipótesis.} Movimiento vertical y empuje ascendente. g constante, atmósfera y suelo omitidos; datos elegidos con empuje inicial superior al peso.

\textbf{Modelo físico.} Ecuación de masa variable con gravedad.
\[m(t)=m_0-\dot m_pt,\quad F=c\dot m_p\]
\[\dot v=c\dot m_p/m(t)-g\]
\[v_b-v_0=c\ln[m_0/(m_0-\dot m_pt_b)]-gt_b\]
\textbf{Desarrollo.}
\begin{enumerate}
\item La masa final es $427-2(55.4)=316.2$ kg, positiva. El empuje es $2450\cdot2=4900$ N.
\item La integral propulsiva vale $736.0023\,\mathrm{m/s}$. La proyección gravitatoria vertical de este caso es $g t_b=543.2884$ m/s, por lo que el incremento neto es $192.7139\,\mathrm{m/s}$.
\item Inicialmente $F/(m_0g)=1.170166>1$. La aceleración inicial es $1.66876\,\mathrm{m/s^2}$ y la final $5.689871\,\mathrm{m/s^2}$. Empuje constante no produce aceleración constante.
\end{enumerate}
\textbf{Operaciones para contrastar el cálculo.}
\begin{enumerate}

\item Magnitud 1.

\[\fitmath{y_{1}=2450\cdot \ln\left(\frac{427}{316.2}\right)-9.80665\cdot 55.4\approx 192.7139\,\mathrm{m/s}}\]

\item Magnitud 2.

\[\fitmath{y_{2}=427-2\cdot 55.4\approx 316.2\,\mathrm{kg}}\]

\end{enumerate}

\textbf{Resultados.} $m_f=316.2$ kg; $\Delta v_{ideal}=736.0023$ m/s; $\Delta v_g=543.2884$ m/s; $v_b-v_0=192.7139$ m/s.

\textbf{Interpretación.} La duración introduce una diferencia respecto a capacidad ideal. Esta velocidad sigue sin demostrar inserción orbital: el movimiento es vertical y el campo se ha simplificado.

\textbf{Comprobación.} Integrar numéricamente $F/(m_0-\dot m t)-g$ debe coincidir con la forma logarítmica. Al hacer $g\to0$ se recupera Tsiolkovski.

\textbf{Límites.} El término gt no es una pérdida gravitatoria universal. No se modela curvatura, rotación, resistencia ni guiado.
\end{studiumworked}
